\chapter{Méthodologie}
\label{ch:methodo}

Ce chapitre décrit la démarche : la conception d'A2 et son modèle de comptage (section~\ref{sec:a2}), les propriétés qui assurent sa correction et sa sûreté (section~\ref{sec:surete}), puis le protocole d'expérience et d'analyse du code produit (section~\ref{sec:protocole}).

\section{Conception d'A2 : comptage, mise en œuvre, hypothèses et limites}\label{sec:a2}
Cette section répond à QR1 et décrit la mise en œuvre (QR2). Les propositions~\ref{prop:borne} à~\ref{prop:real} sont des énoncés mathématiques \emph{dans un modèle} ; on n'en tire aucune conclusion sur la vitesse (section~\ref{sec:perfborne}).

\subsection{Modèle de comptage}
\begin{definition}[Données]
Un \emph{$m$-vec} est un vecteur $v\in\mathbb{F}_{16}^m$, avec $m$ pair, rangé sur $b=m/2$ octets : $v_{2i}$ est le quartet bas de l'octet $i$ et $v_{2i+1}$ le quartet haut. Les $K$ nombres $c_0,\dots,c_{K-1}$ sont regroupés en $P=\lceil K/2\rceil$ paires (complétées par 0 si $K$ est impair). Pour chaque $k<K$, on calcule $\mathrm{acc}_k\leftarrow\mathrm{acc}_k\oplus c_k v$.
\end{definition}
Pour la paire $p$, on pose $T_p[x]=(c_{2p}x)\oplus((c_{2p+1}x)\ll4)$. Un \vpshufb{} de largeur $W\in\{16,32\}$ octets fait $W$ recherches indépendantes. Chacune donne, pour \emph{un} élément de $v$, les deux produits de la paire.

\begin{definition}[Classe $\mathcal{C}$]
Algorithmes qui (i) obtiennent chaque produit $c_kv_j$ par une recherche dans une table de 16 octets ; (ii) donnent à chaque recherche un quartet seul dans un octet (bit~7 nul) ; (iii) ne font pas deux recherches pour le même couple (élément, paire), sauf pour le recouvrement entre registres ; (iv) utilisent librement XOR, ET, décalages, permutations et chargements. Sont \emph{hors} de $\mathcal{C}$ : GFNI, les représentations \emph{bitsliced}, les tables de 256 octets.
\end{definition}
La quantité comptée est $N(\text{algo},m)$ : le nombre de \vpshufb{} exécutés \emph{par paire de nombres et par $m$-vec}.

\subsection{Énoncés}
\begin{proposition}[Limite]\label{prop:borne}
Tout algorithme de $\mathcal{C}$ vérifie $N\ge\lceil m/32\rceil$.
\end{proposition}
\begin{proof}
Chaque élément $v_j$ ($j<m$) demande au moins une recherche par paire, et une instruction en fait au plus $W\le32$. Donc $32N\ge m$.
\end{proof}

\begin{proposition}[Schéma officiel avec recouvrement]\label{prop:off}
Le schéma qui couvre les $b$ octets avec $\lceil b/32\rceil$ chargements de 32 octets (le dernier est reculé) et qui traite chaque registre avec deux \vpshufb{} (impairs, pairs) a $N_{\mathrm{off}}(m)=2\lceil m/64\rceil$.
\end{proposition}
\begin{proof}
Il y a $\lceil b/32\rceil=\lceil m/64\rceil$ chargements. Chacun donne 32 indices impairs et 32 indices pairs, donc deux \vpshufb{}.
\end{proof}

\begin{proposition}[Mise en œuvre de la limite si $0<r\le32$]\label{prop:real}
Écrivons $m=64q+r$ avec $0\le r<64$, et $t=r/2$ le nombre d'octets de la fin du vecteur. Si $0<r\le32$, l'algorithme suivant est dans $\mathcal{C}$ et a $N=2q+1=\lceil m/32\rceil$ :
\begin{enumerate}[nosep]
\item pour $g<q$ : charger l'intervalle $[32g,32g+32)$ et le traiter avec deux \vpshufb{} ymm ;
\item pour la fin du vecteur : charger les $t$ octets de $[32q,32q+t)$, construire $X=[\text{quartets bas des }t\text{ octets}\mid\text{quartets hauts des }t\text{ octets}]$ (cela met les $r\le32$ indices dans un seul registre), puis appliquer \emph{un seul} \vpshufb{} ;
\item remettre en ordre les sorties de la fin du vecteur (même formule que pour un registre complet, limitée aux $r$ positions).
\end{enumerate}
\end{proposition}
\begin{proof}
On compte $2q+1$ instructions, et $\lceil m/32\rceil=2q+\lceil r/32\rceil=2q+1$ quand $0<r\le32$.
\end{proof}
\emph{Conditions de mise en œuvre} (le code doit les respecter ; le modèle ne les garantit pas) : (a) le chargement de la fin du vecteur ne lit que des octets de l'entrée ; (b) si $r\le16$ ($t\le8$), $X$ tient dans un registre de 128 bits et on utilise la \emph{voie basse} de $T_p$ ; si $16<r\le32$, $X$ occupe un registre de 256 bits (voie basse = impairs, voie haute = pairs), ce qui demande que les deux voies soient égales ; (c) le coût de préparation de la fin du vecteur est payé une fois par entrée, et celui de la remise en ordre une fois par paire.

\begin{corollaire}[Limite exacte]
Pour $m\le192$, $\Delta(m):=N_{\mathrm{off}}(m)-\lceil m/32\rceil=1$ si $0<(m\bmod64)\le32$, et $0$ sinon. Dans le second cas, le code officiel atteint déjà la limite : aucun algorithme de $\mathcal{C}$ ne peut faire mieux en nombre de \vpshufb{}.
\end{corollaire}

\begin{table}[ht]\centering\small
\caption{Application aux paramètres actuels.}\label{tab:delta}
\begin{tabular}{@{}lrrrrrrl@{}}\toprule
 & $m$ & $q$ & $r$ & $N_{\mathrm{off}}$ & $\lceil m/32\rceil$ & $\Delta$ & Statut\\\midrule
MAYO-1 & 80 & 1 & 16 & 4 & 3 & 1 & implémentée (A2), fin de 8 octets\\
MAYO-2 & 64 & 1 & 0 & 2 & 2 & 0 & sans objet\\
MAYO-3 & 108 & 1 & 44 & 4 & 4 & 0 & impossible dans $\mathcal{C}$\\
MAYO-5 & 142 & 2 & 14 & 6 & 5 & 1 & implémentée (A2), fin de 7 octets + 1 de bourrage\\\bottomrule
\end{tabular}
\end{table}

\paragraph{Ce qui est fait et ce qui ne l'est pas.} Le prototype traite les fins de vecteur d'au plus 16 éléments ($t\le8$, $r\le16$). La variante pour $16<r\le32$ \textbf{n'est ni réalisée ni testée} : ce n'est pas un résultat acquis. Aucun des quatre paramètres actuels n'a $r\in(16,32]$.

\subsection{Conception d'A2}
\subsubsection{Structure du noyau}
Pour une entrée de $P_1$ et une ligne d'accumulateurs de $K$ vecteurs, A2 garde tels quels les $q=\lfloor m/64\rfloor$ registres ymm complets sur les $32q$ premiers octets. Il traite ensuite la fin du vecteur ainsi :
\begin{enumerate}[nosep]
\item \textbf{Préparation (une fois par entrée)} : chargement de 8 octets (\code{_mm_loadl_epi64}), décalage de 4 bits (\code{vpsrlw}), entrelacement \code{vpunpcklqdq} (quartets bas dans les octets 0 à 7, quartets hauts dans les octets 8 à 15), ET avec \code{0x0f}. Cela fait quatre instructions (cinq dans la variante SAFE, qui ajoute un décalage d'octets).
\item \textbf{Produits} : un \code{vpshufb} xmm avec la \emph{moitié basse} de la table de multiplication déjà chargée, et un \code{vpxor} xmm dans un accumulateur xmm.
\item \textbf{Remise en ordre (fin de ligne)} : séparation impair/pair comme pour la partie ymm, puis lecture, modification et écriture d'\emph{un seul} limbe de 64 bits (le limbe d'indice $4q=L-1$, soit 4 pour MAYO-1 et 8 pour MAYO-5).
\end{enumerate}
Dans l'assembleur produit par GCC~16.2.0, la boucle interne passe donc, par entrée, de 20 à 15 \code{vpshufb} pour MAYO-1 (10 ymm + 5 xmm) et de 36 à 30 pour MAYO-5 (24 ymm + 6 xmm).

\subsubsection{Variante par défaut et variante SAFE (MAYO-5)}
Pour MAYO-1, $b=40=32+8$ : la fin du vecteur est exactement la fin de l'entrée, $[32,40)$. Pour MAYO-5, $b=71$ et $64+8=72$ :
\begin{itemize}[nosep]
\item \emph{par défaut} : lit $[64,72)$, donc \textbf{un octet après} l'entrée (comme le noyau officiel, dont le dernier chargement couvre $[40,72)$) ;
\item \emph{SAFE} : lit $[63,71)$ puis décale le registre de 8 bits vers les poids faibles ; l'octet 63 est éliminé, les octets 64 à 70 passent aux positions 0 à 6, et la position 7 reçoit 0. \textbf{Toutes les lectures restent dans l'entrée.}
\end{itemize}
Si l'on intègre un jour A2, on ne retiendrait que la variante SAFE. Les quartets de bourrage (positions 142 et 143 de MAYO-5, car $16L=144$) n'ont pas d'effet sur les positions utiles (section~\ref{sec:bourrage}).

\subsubsection{Préconditions}
(P1) $P_1$ est rangé à plat, sans bourrage entre les entrées, à l'adresse $P_1+e\cdot b$ ; (P2) les accumulateurs ont $L$ limbes, avec du bourrage de quartets seulement à la fin du dernier limbe ; (P3) les tables de multiplication sont construites par la routine officielle (voies égales, lemme~\ref{lem:voies}) ; (P4) $m\in\{80,142\}$ : les noyaux A2 ne sont pas actifs pour les autres paramètres (condition de compilation \code{MAYO_TAIL_XMM_ACTIVE}) ; (P5) compilation sans \code{ENABLE_PARAMS_DYNAMIC}.

\subsection{La limite n'implique pas l'accélération}\label{sec:perfborne}
Faire moins de \vpshufb{} ne garantit pas d'aller plus vite : il faut ajouter le coût de préparation de la fin du vecteur, les dépendances entre instructions, les rechargements depuis la pile, les chargements non alignés et le débit des unités de calcul. Si seul le débit du \vpshufb{} (un par cycle sur Skylake~\cite{nibbling}) limitait la boucle, le rapport de temps attendu serait $N_{A2}/N_{\mathrm{off}}$ : $15/20=0{,}750$ (MAYO-1) et $30/36=0{,}833$ (MAYO-5).

\begin{table}[ht]\centering\small
\caption{Rapport attendu par le seul débit de \vpshufb{} et rapport mesuré (S1, A2/officiel, noyaux isolés).}\label{tab:borne-mesure}
\begin{tabular}{@{}lccl@{}}\toprule
Noyau & Limite « débit seul » & Mesuré (S1) & Lecture\\\midrule
MAYO-1 \code{P1_times_Vt} & 0,750 & 0,780 & plus lent que la limite\\
MAYO-1 \code{P1t_times_Vt} & 0,750 & 0,710 & plus rapide que la limite\\
MAYO-5 \code{P1_times_Vt} & 0,833 & 0,807 & plus rapide que la limite\\
MAYO-5 \code{P1t_times_Vt} & 0,833 & 0,880 & plus lent que la limite\\\bottomrule
\end{tabular}
\end{table}
Deux cas sur quatre sont plus rapides que ce que permettrait le seul comptage des \vpshufb{} : d'autres effets jouent (par exemple moins de rechargements depuis la pile : de 11 à 5 pour MAYO-1 \code{P1_times_Vt}, de 30 à 24 pour MAYO-5), mais ce n'est pas établi. Le modèle fondé sur le débit n'explique donc pas à lui seul les gains. Une hypothèse de performance H-P (« sur un processeur où le débit des permutations est le facteur limitant, le temps est au moins $N\cdot P$ cycles par entrée ») ne peut se tester qu'avec des cycles d'horloge, pas avec des ticks TSC. Elle vient après ce travail et reste ouverte.

\section{Propriétés de correction et de sûreté}\label{sec:surete}
Cette section répond à la partie « sûreté » de QR2. Chaque propriété porte une étiquette : \Pc{} (déduite du code et de ses invariants, par relecture et par des assertions à la compilation), \Pe{} (vérifiée par des tests) ou \Hy{} (hypothèse sur la machine ou le compilateur). Une propriété \Pe{} seule n'est pas une démonstration. Notation : $n_y=\lfloor m/64\rfloor$ (1 pour MAYO-1, 2 pour MAYO-5), $T_{\mathrm{OFF}}=32n_y$ (32 ; 64), $T_{\mathrm{LIMB}}=4n_y$ (4 ; 8).

\subsection{Limites de lecture de $P_1$}
L'entrée $e$ occupe les octets $[eb,(e+1)b)$.
\begin{itemize}
\item \Pc{} Les $n_y$ chargements ymm lisent $[eb+32g,eb+32g+32)$ pour $g<n_y$ ; comme $32n_y\le b$ ($32\le40$ ; $64\le71$), ils restent dans l'entrée.
\item \Pc{} MAYO-1 : la fin du vecteur lit $[eb+32,eb+40)$, exactement la fin de l'entrée.
\item \Pc{} MAYO-5, SAFE : la fin du vecteur lit $[eb+63,eb+71)$ puis décale de 8 bits ; toutes les lectures restent dans l'entrée.
\item \Pc{} (variante par défaut, non retenue) MAYO-5 lit $[eb+64,eb+72)$ : un octet après l'entrée. Pour $e<N_{\mathrm{ent}}-1$, c'est le premier octet de l'entrée suivante. Pour la dernière entrée, c'est l'octet 720\,863 du tampon, qui est alloué (zone de $8\cdot$\code{P1_LIMBS_MAX}$=731\,016$ octets). Des \code{_Static_assert} imposent $N_{\mathrm{ent}}b+\text{OVERREAD}\le8\cdot$\code{P1_LIMBS_MAX} et OVERREAD$\le1$.
\item \Pe{} Test « garde-page » (la fin de $P_1$ est collée à une page interdite, sans marge) : le code officiel et A2 par défaut provoquent une erreur SIGSEGV pour MAYO-5, A2 SAFE non ; pour MAYO-1, aucun des trois. Ce test contrôle la propriété ci-dessus sans la démontrer.
\end{itemize}
Le noyau officiel de MAYO-5 lit lui aussi 1 octet de trop (son dernier ymm couvre les octets 40 à 71), et celui de MAYO-3 en lit 2 (son second ymm couvre $[24,56)$ pour $b=54$). A2 SAFE n'a besoin d'aucune marge.

\subsection{Limites d'écriture}
Pour la ligne $r$, l'accumulateur est à $\mathrm{acc}+rKL$ (en limbes) ; le vecteur $k$ occupe $[kL,(k+1)L)$.
\begin{itemize}
\item \Pc{} Les écritures ymm touchent les limbes $[kL+4g,kL+4g+4)$ pour $g<n_y$, et jamais au-delà du limbe $4n_y-1\le L-1$.
\item \Pc{} La fin du vecteur est traitée par lecture, modification puis écriture d'un seul limbe de 64 bits, d'indice $T_{\mathrm{LIMB}}=L-1$, pour $k<K$ (et $k+1<K$ pour le second, avec une garde explicite). Toutes les écritures sont donc dans les $VKL$ limbes de \code{acc}. Le noyau officiel écrit un ymm à $L-4$ : même étendue.
\item \Pe{} Un « canari » placé après \code{acc} n'est pas modifié ; $P_1$, passé en \code{const}, n'est pas modifié non plus (la copie de l'entrée est inchangée).
\end{itemize}

\subsection{Invariant des deux voies des tables}
A2 lit la \emph{voie basse} de chaque table et suppose qu'elle est égale à la voie haute.
\begin{lemme}\label{lem:voies}\Pc{} Pour tout $i$, \code{V_multabs[i]} a ses deux voies de 128 bits égales.
\end{lemme}
\emph{Justification.} (i) Chacune des quatre lignes de 32 octets de \code{__gf16_mulbase} répète deux fois la même suite de 16 octets (constantes lues dans \code{arithmetic_common.h}, lignes 12 à 17). (ii) \code{tbl32_gf16_multab2(b)} est un OU exclusif de ET entre ces lignes et des masques obtenus par \code{_mm256_set1_epi16(b & 0xf)} décalé puis comparé ; ces opérations agissent par élément de 16 bits, donc voie par voie, avec des masques identiques dans les deux voies. (iii) \code{mayo_V_multabs} combine deux telles tables par \code{^} et \code{_mm256_slli_epi16(.,4)}, aussi voie par voie. Par récurrence, l'égalité est conservée.

\Pe{} Nous avons vérifié 272 tables avec le mode \code{MAYO_TAIL_XMM_CHECK} (arrêt en cas d'écart), sur les tables données aux noyaux. C'est une \emph{condition préalable} d'A2 : si les tables étaient construites autrement (autre routine, autre version du code), il faudrait revérifier cet invariant.

\subsection{Influence des quartets de bourrage}\label{sec:bourrage}
Seul MAYO-5 a des quartets de bourrage (positions 142 et 143 ; pour MAYO-1, $80=5\cdot16$).
\begin{itemize}
\item \Pc{} \emph{Chaque position est indépendante.} La sortie à la position $j$ ne dépend que des entrées à la position $j$ de chaque $m$-vec : une recherche lit un index et une table, jamais un autre index. Les décalages \code{srli/slli_epi16} mélangent deux octets voisins, mais chaque résultat est suivi d'un ET avec \code{0x0f} (ou vient d'un registre déjà masqué).
\item \Pc{} Les deux seules différences entre le code officiel et A2 aux positions de bourrage sont : pour l'officiel, les produits de l'octet 71 (premier octet de l'entrée suivante) ; pour A2 SAFE, 0.
\item \Pc{} Les noyaux alimentent \code{Pv}, puis \code{Vt_times_Pv} produit \code{VP1V}, qui est utilisé par \code{compute_rhs}. C'est le premier endroit qui mélange des positions, et il met à zéro les quartets de bourrage avant de les utiliser (\code{src/mayo.c}, lignes 50 à 60, masque $(1\ll((m\bmod16)\cdot4))-1$). Entre les noyaux et \code{compute_rhs}, toutes les opérations travaillent position par position.
\item \Pe{} Tests unitaires (20\,000 cas par taille, 8 familles d'entrées) : 0 écart sur les $M/2$ octets utiles et sur les limbes complets après masquage. Sur 30\,010 signatures appariées de MAYO-5 (15\,005 pour MAYO-1 et pour MAYO-2), aucune signature n'est différente.
\end{itemize}
\textbf{Limite.} Cette propriété suppose que \code{Pv} n'est utilisé que par les chemins listés, ce que nous avons établi en cherchant dans le code source de la version AVX2. Si une modification future lisait le bourrage de \code{Pv}, il faudrait réexaminer la question.

\subsection{Temps constant : ce qui est établi et ce qui ne l'est pas}
\begin{itemize}
\item \Pc{} Les bornes des boucles ($V$, $K$) et les adresses de chargement ($eb$) ne dépendent que de paramètres publics. Les tables calculées à partir du vinaigre ne sont que des données pour \vpshufb{} : ce ne sont jamais des adresses ni des conditions. Le code officiel a la même structure.
\item \Pc{}$\to$à vérifier sur le programme compilé : à la remise en ordre de la fin du vecteur, A2 ajoute deux \code{vmovq} d'un registre vectoriel vers un registre général (\code{ar[...] ^= (uint64_t)_mm_cvtsi128_si64(...)}). Des données liées au secret passent donc par un registre général, mais seulement comme donnée d'un XOR vers la mémoire. Nous avons inspecté l'assembleur (script \code{inspect_asm.py}) et n'avons trouvé aucun usage non conforme (adresse, branche, drapeau). Ce résultat vaut pour \emph{une} version de GCC, pas en général.
\item \Pe{} (indice) Valgrind en mode temps constant (données secrètes marquées comme indéfinies), dans la VM seulement : il repère les branches et les adresses qui dépendent de données marquées, sur les exécutions testées.
\item \Hy{} Les instructions \vpshufb{}, \code{vpxor}, \code{vpand}, les décalages, les chargements et les écritures prennent un temps qui ne dépend pas des données sur les processeurs visés (d'après la documentation du constructeur, non vérifiée ici).
\end{itemize}
\textbf{Nous ne fournissons aucune preuve formelle de temps constant.} Que les KAT réussissent et qu'il n'y ait pas de différence de résultat ne démontre pas le temps constant.

\subsection{Correction fonctionnelle}
Les KAT officiels réussissent pour MAYO-1, 2, 3 et 5 (toutes les sessions), ainsi que les tests du schéma et les tests unitaires précédents. Il y a zéro signature différente dans toutes les paires A/B, et le nombre de tentatives de signature est le même. La correction est donc établie \emph{sur les entrées testées} ; l'argument général repose sur les propriétés \Pc{} ci-dessus.

\section{Protocole expérimental}\label{sec:protocole}
L'objectif est qu'un lecteur qui a le dépôt, les archives et une machine comparable puisse refaire chaque analyse et, au bruit près, chaque mesure. Les commandes sont à l'annexe~\ref{app:inventaire}.

\subsection{Plateforme}
\paragraph{Machine physique (toutes les mesures citées au chapitre~\ref{ch:resultats}).} Portable Intel Core i9-9880H (8 cœurs, 16 threads, microcode \code{0x104}) ; \code{systemd-detect-virt} = \code{none}, drapeau hyperviseur = 0 ; Kali 7.1.5 (démarrage sur disque, double amorçage) ; GCC 16.2.0 (Debian 16.2.0-3), cmake 4.3.4, Python 3.14.8, numpy 2.4.6. Gouverneur \code{performance}, turbo actif, SMT actif, secteur branché. Processus épinglé sur le cœur 1 (\code{taskset}). La fréquence instantanée du cœur n'est enregistrée qu'au lancement (de 800 MHz au repos à 4\,540 MHz selon les sessions) ; elle n'est pas une fréquence soutenue. Horloge du portable à UTC+2 ; les journaux \code{machine.txt} sont en UTC.

\paragraph{Environnement virtuel (historique, non utilisé pour les résultats de performance).} VM KVM à 2 cœurs, GCC 13.3.0, cmake 3.28.3. Les mesures de cette VM (qui ont orienté les décisions de conception) ne sont jamais présentées comme performances matérielles.

\paragraph{Compilation.} Options de la bibliothèque : \code{-O3 -funroll-loops -march=haswell -maes}, identiques pour les deux versions comparées. Pour l'analyse statique, la commande exacte de chaque objet est lue dans \code{compile_commands.json}.

\subsection{Code comparé}
Base amont : dépôt officiel \code{PQCMayo/MAYO-C}, commit \code{74de637a7c4a7f9fafa888183bc05f5cf920cd2d}. Les modifications sont appliquées par \code{git am} dans l'ordre : \code{0001} (prototype A2), \code{0002} (sondes de décomposition, inertes par défaut), \code{0003} (contrôle de l'effet des sondes), \code{0004} (analyse statique E1, aucun noyau modifié). Les identifiants de commit changent avec \code{git am} ; les arbres, non (annexe~\ref{app:inventaire}). L'arbre des mesures S1/S2 est celui après \code{0001} ; celui de PH1/CT1 après \code{0003}. Les noyaux optimisés n'ont pas été modifiés pendant la phase d'analyse.

\subsection{Protocole de mesure apparié}
Pour chaque niveau mesuré (\code{P1_times_Vt}, \code{P1t_times_Vt}, \code{m_calculate}, \code{mayo_sign}), un processus exécute des paires (officiel, A2) dans un ordre \emph{aléatoire}, avec la même graine et le même message pour les deux versions. Le temps est lu par \code{rdtscp} (ticks TSC). L'ordre d'exécution de chaque paire est enregistré (\code{xmm_first}). Un processus comprend 3\,001 paires. Les signatures des deux versions sont comparées octet par octet, et le nombre de tentatives de signature est vérifié.

\subsection{Unité statistique et intervalles}
Les 3\,001 paires d'un même processus ne sont pas des répétitions indépendantes (même état de la machine, même fréquence). L'intervalle à 95\,\% « naïf », calculé sur les paires, est d'environ $\pm0{,}0001$ pour \code{mayo_sign}, alors que la dispersion entre processus est 5 à 15 fois plus grande. \textbf{L'unité statistique est le processus.} Pour chaque niveau et chaque processus, on calcule la médiane du rapport A2/officiel. La valeur donnée est la médiane de ces médianes, avec un intervalle à 95\,\% obtenu par \emph{bootstrap sur les processus}. Nous testons une seule différence entre sessions (S2$-$S1) : avec deux démarrages, on ne peut pas estimer à part la variation « entre sessions » ; nous donnons donc l'écart et son intervalle bootstrap.

\paragraph{Effet d'ordre.} Pour chaque processus, on calcule la médiane du rapport sur les paires « A2 en premier » et sur les paires « officiel en premier ». L'effet d'ordre est la différence, avec un intervalle bootstrap sur les processus ; on le juge par rapport au gain.

\subsection{Sessions}
\begin{table}[ht]\centering\small
\caption{Sessions de mesure sur machine physique (dates UTC).}\label{tab:sessions}
\begin{tabularx}{\linewidth}{@{}lXXX@{}}\toprule
Session & Date, démarrage & Plan & Objet\\\midrule
S1 & 19 oct. 21:58--22:43 ; un démarrage & 4 campagnes $\times$ 5 processus $\times$ 3\,001 paires (graines 1--20) & gains\\
S2 & 20 oct. 12:10 ; nouveau démarrage & 5 processus (graines 21--25) & reproductibilité\\
PH1 & 21 oct. 19:13 ; nouveau démarrage & MAYO-5 : 10 processus ; MAYO-1, MAYO-2 : 5 & décomposition instrumentée\\
CT1 & 21 oct. 19:28 ; même démarrage que PH1 & 5 configurations $\times$ 10 processus (graines 41--50) & effet des sondes\\\bottomrule
\end{tabularx}
\end{table}
Les scripts de lancement refusent d'écraser un dossier ou une archive existants : aucune campagne n'écrase la précédente. S1 et S2 comparent deux démarrages différents avec le même code, la même machine et le même compilateur. Pour les temps absolus, S2 est environ 7\,\% plus lent que S1 pour tous les niveaux, témoins compris ; nous n'avons pas mesuré la cause, et les rapports ne sont pas touchés.

\subsection{Instrumentation de décomposition}
Des sondes (\code{rdtscp} précédé d'un \code{lfence}) mesurent, pour chaque signature et chaque version, le temps total, le nombre de tentatives et, pour chaque phase, le temps et le nombre d'appels. Elles sont inactives par défaut : sans la macro correspondante, le code objet de \code{mayo.c} et de \code{arithmetic.c} est identique à l'original (vérifié par \code{cmp} et \code{objdump}). Le patch \code{0003} ajoute un masque pour activer chaque sonde séparément. Les cinq configurations de CT1 sont : \code{plain} (bibliothèque d'origine, aucune sonde), \code{none} (sondes compilées mais toutes désactivées), \code{esk} (sonde sur \code{expand_sk_internal} seule), \code{mcalc} (sonde sur \code{m_calculate_M_and_VPV} seule) et \code{all} (toutes les sondes).

\subsection{Analyse statique du code généré (E1)}
Pour chaque construction, nous recompilons l'objet avec la commande enregistrée, nous le désassemblons (\code{objdump -d -M intel}) et nous repérons les boucles de travail par leur nombre de \code{vpshufb} (36 ymm pour l'officiel de MAYO-5 ; 24 ymm + 6 xmm pour A2). Les comptes sont ramenés au facteur de déroulage de boucle (ici $u=1$). Nous identifions le noyau par le \emph{pas d'adresse sur $P_1$} : \code{add reg,0x47} (entrées de 71 octets, le pas d'une entrée) pour \code{P1_times_Vt} et \code{imul ...,0x47} pour \code{P1t_times_Vt}. Les repères de débogage de GCC ne sont pas fiables (les corps identiques sont fusionnés et étiquetés \code{P1_times_O}). Pour attribuer une différence de code à une cause, il faut d'abord avoir écarté les autres modifications du code (section~\ref{sec:e1b}).

\subsection{Protocole E1b (analyse statique de l'origine des écritures)}\label{sec:e1b}
E1b recompile \code{src/mayo.c} avec \emph{la commande exacte} de E1 plus une seule option en plus, sans aucune mesure de temps. Constructions : \code{P_none} (sans sondes), \code{P_all} (toutes les sondes), \code{TAB} (construction des mesures S1/S2). Variantes : \code{base} (contrôle : doit redonner exactement le code de E1) ; \code{dump} (fichiers de sortie \code{-fdump-tree-optimized}, \code{-fdump-tree-sra-details}, \code{-fdump-tree-esra-details}, \code{-fdump-rtl-expand}, \code{-fdump-rtl-ira}, \code{-fopt-info-inline-optimized} ; contrôle : le code doit être identique) ; \code{nosra} (\code{-fno-tree-sra}) ; \code{noipasra} (\code{-fno-ipa-sra}) ; \code{nodse} ; \code{nopre} ; \code{nofre}.

Pour chaque variante, le rapport donne : les boucles de travail avec leur nombre d'instructions et d'écritures de 32 octets ; le nombre d'instructions de \code{sign_signature_internal} \emph{en dehors des boucles de travail} ; la distance entre les histogrammes de mnémoniques, hors boucles, par rapport à \code{base} ; et le nombre de lignes GIMPLE qui mentionnent \code{temp}. \textbf{Règle de lecture :} on ne lit jamais seule la disparition des écritures dans une variante. On ne l'attribue à une passe que si la distance hors boucles est faible, c'est-à-dire si rien d'autre n'a changé.


\begin{table}[ht]\centering\small
\caption{Issues possibles d'E1b et conclusions autorisées (fixées avant exécution).}\label{tab:e1b-issues}
\begin{tabularx}{\linewidth}{@{}lXX@{}}\toprule
 & Observation (GCC~16.2.0) & Conclusion autorisée / non autorisée\\\midrule
1 & Doublon absent avec \code{-fno-tree-sra} \emph{et} distance en dehors des boucles faible & la passe conditionne le doublon dans cette compilation / non : que \code{temp[]} en soit la cause sans les fichiers d'analyse\\
2 & Doublon absent mais distance élevée & l'attribution est \emph{non établie}\\
3 & Doublon présent dans toutes les variantes & les passes testées ne le suppriment pas, hypothèse non confirmée / non : que le doublon soit sans cause compilateur\\
4 & \code{base} $\ne$ code de E1 & E1 non reproduit : arrêter l'interprétation\\\bottomrule
\end{tabularx}
\end{table}

\subsection{Traçabilité et reproduction}
Les données brutes (S1, S2, PH1, CT1) sont publiées avec leurs empreintes SHA-256 (annexe~\ref{app:inventaire}). Le script de reproduction réextrait les archives, vérifie les empreintes, relance les neuf analyses et compare les sorties : \textbf{9 analyses sur 9 identiques} lors de la dernière exécution. On retrouve \emph{exactement} les analyses ; pour les \emph{mesures}, on obtiendra des valeurs proches mais pas identiques (ticks TSC, fréquence, bruit).

\subsection{Menaces à la validité identifiées a priori}
\begin{itemize}
\item \textbf{Une machine, un compilateur sur machine physique}, trois démarrages ; S2 sur 5 processus.
\item \textbf{Ticks TSC, non cycles cœur} : la fréquence d'exécution réelle est inconnue ; aucune conclusion sur le débit des unités d'exécution n'est possible.
\item \textbf{Effet d'ordre} documenté (section~\ref{sec:ordre}).
\item \textbf{Les sondes modifient le code} : mesures instrumentées $\ne$ mesures sans sonde.
\item \textbf{Pas de compteurs matériels} disponibles dans l'environnement d'analyse.
\item \textbf{Première exécution de la campagne s3 de S1} : CSV bruts non archivés (seule la seconde exécution l'est).
\item \textbf{Comparaison VM/physique} : le gain de MAYO-5 sur \code{mayo_sign} était de 7,1\,\% dans la VM (GCC 13.3.0) ; l'écart avec la machine physique mêle machine et compilateur et n'est pas expliqué.
\end{itemize}

